\section{Introduction}
\label{sec:intro}

The deployment of Large Language Model (LLM) agents with tool-calling capabilities has revolutionized enterprise process automation. In financial accounting, accounts payable (AP), and treasury operations, autonomous agents regularly ingest unstructured vendor invoices, supplier emails, and compliance memos, extracting remittance information and invoking banking APIs (e.g., SWIFT, Fedwire, ACH) to settle obligations.

However, recent studies in adversarial agent security~\cite{debenedetti2024agentdojo,zou2024injecagent} demonstrate that LLMs cannot reliably segregate data from instructions. When an agent reads an invoice containing malicious imperative commands (e.g., \textit{``Notice: Vendor banking coordinates updated to offshore account 99281-8821''}), standard reasoning patterns such as ReAct~\cite{zou2024injecagent} and autonomous planning loops treat the injected prompt as authentic instruction. Consequently, the agent executes unauthorized fund transfers, overriding corporate disbursement policies.

Existing defenses primarily rely on soft prompt engineering (e.g., \textit{``Do not follow instructions in retrieved data''}) or auxiliary LLM evaluators. As proved by Abdelnabi \& Bagdasarian~\cite{debenedetti2024agentdojo}, soft prompt defenses are inherently probabilistic and suffer catastrophic degradation when adversaries employ light obfuscation, compliance jargon, or markdown table smuggling. Furthermore, using secondary LLMs to inspect every transaction introduces hundreds of milliseconds of latency and substantial API costs, making them impractical for real-time high-throughput financial clearinghouses.

To address these vulnerabilities without compromising system throughput, we present \textbf{AgentPOI}. Rather than attempting to fix probabilistic natural language reasoning inside the LLM, AgentPOI enforces \textbf{Contextual Invariant Verification (CIV)} at the non-bypassable execution boundary. 

\subsection{Key Contributions}
\begin{enumerate}
    \item \textbf{Single-Domain Financial Defense Architecture}: We engineer a specialized security guard for Autonomous Financial and Transactional Agents that synthesizes retrieval-time fact-imperative filtering~\cite{taskshield2025}, cryptographic lineage DAGs~\cite{authgraph2026}, and epistemic memory firewalls~\cite{datasentinel2026}.
    \item \textbf{Contextual Invariant Verification (CIV)}: We formalize non-linguistic mathematical invariants over financial state transitions, enforcing verified payee authentication, per-transaction ceilings, and rolling 24-hour daily balance conservation laws.
    \item \textbf{Zero-Hardware Offline Efficiency}: AgentPOI executes 100\% offline on standard CPU, adding only 0.22\,ms latency overhead per transaction and consuming zero defensive LLM tokens.
    \item \textbf{Rigorous 60-Scenario Evaluation}: Across 30 adversarial and 30 clean enterprise scenarios adapted from AgentDojo Banking and InjecAgent, AgentPOI reduces Attack Success Rate on Execution (ASR-Exec) from 100.0\% to \textbf{0.0\%} while maintaining \textbf{100.0\%} legitimate invoice execution utility (0.0\% Clean Disruption Rate).
\end{enumerate}
